% fig7_28.tex — 7-42: 驻波（绳左端固定在墙上，O点距墙5λ/4）
\begin{tikzpicture}[font=\small,>=stealth]
  % 墙
  \draw[thick] (-0.3,1.5) -- (-0.3,-1.5);
  \foreach \y in {-1.3,-1.1,...,1.3}
    \draw (-0.3,\y) -- (-0.6,{\y+0.15});
  \node[left] at (-0.3,0) {墙};
  % 绳子（平衡位置为水平线）
  \draw[thick] (-0.3,0) -- (8,0);
  % 驻波波形（包络线）
  % 5λ/4 = 1.25λ，从墙到O点
  % 节点在墙处（固定端），在 λ/4, 3λ/4, 5λ/4 处有节点
  % 设λ = 2单位，则5λ/4 = 2.5单位
  % 绘制驻波的两个极端包络
  \draw[thick,blue!60,domain=0:2.5,samples=100] plot (\x,{1.2*sin(deg(pi*\x/1.25))});
  \draw[thick,blue!60,domain=0:2.5,samples=100] plot (\x,{-1.2*sin(deg(pi*\x/1.25))});
  % 虚线填充中间区域
  \draw[blue!20,fill=blue!10,domain=0:2.5,samples=100] plot (\x,{1.2*sin(deg(pi*\x/1.25))})
    -- (2.5,0) -- cycle;
  % O点标记
  \fill (2.5,0) circle (2pt) node[below] {$O$};
  % 节点标注
  \foreach \x in {0,1.25,2.5}
    \fill (\x,0) circle (1.5pt);
  \node[below,font=\footnotesize] at (0.625,0) {波腹};
  \node[below,font=\footnotesize] at (1.875,0) {波腹};
  \node[above,font=\footnotesize] at (1.25,0.15) {节点};
  % λ/4 标注
  \draw[<->] (0,-1.8) -- (0.625,-1.8) node[midway,below,font=\footnotesize] {$\lambda/4$};
  % 5λ/4 标注
  \draw[<->] (0,-2.2) -- (2.5,-2.2) node[midway,below,font=\footnotesize] {$\frac{5}{4}\lambda$};
  % 墙到O标注
  \draw[|<->|] (-0.3,-2.6) -- (2.5,-2.6) node[midway,below] {墙到 $O$};
\end{tikzpicture}
